% !PRINTED 541
\weberpage{541}

%<<CONT
$F(x, \eta) = 0$ ist aber wieder cyklisch in $\Omega(\eta)$, da die cyklischen Functionen ihrer Wurzeln diesem Körper angehören.

Die cyklischen Gleichungen haben, wie alle \gesperrt{Abel}'schen Gleichungen, die Eigenschaft, dass jede Wurzel rational durch jede andere ausdrückbar ist. Hier lassen sich diese Ausdrücke folgendermaassen cyklisch anordnen. Sind $\alpha, \alpha_1, \alpha_2 \ldots \alpha_{m-1}$ die Wurzeln der cyklischen Gleichung $F(x) = 0$, so ist die ganze Function $(m-1)^{\text{ten}}$ Grades von $x$
\[
F(x)\left(\frac{\alpha_1}{x-\alpha} + \frac{\alpha_2}{x-\alpha_1} + \cdots + \frac{\alpha}{x-\alpha_{m-1}}\right) = \Psi(x)
\]
ungeändert durch die Permutation $\pi$ und also in $\Omega$ enthalten. Wenn wir darin $x = \alpha, \alpha_1 \ldots \alpha_{m-1}$ setzen und das Zeichen
\[
\frac{\Psi(x)}{F'(x)} = \Theta(x)
\]
einführen, so folgt
\weq{11}{\alpha_1 = \Theta(\alpha), \alpha_2 = \Theta(\alpha_1), \ldots \alpha_{m-1} = \Theta(\alpha_{m-2}), \alpha = \Theta(\alpha_{m-1}),}
und dies gilt, mag $F(x)$ reducibel oder irreducibel sein, wenn nur $F(x)$ und $F'(x)$ keinen gemeinsamen Theiler haben.

Die Auflösung der cyklischen Gleichungen $m^{\text{ten}}$ Grades ist hierdurch abhängig gemacht von der Lösung cyklischer Gleichungen, deren Grade die Primfactoren von $m$ sind. Ist also z.\,B. $m$ eine Potenz von $2$, so wird die Lösung durch eine Reihe von Quadratwurzeln bewerkstelligt. Auf diesem Wege hat \gesperrt{Gauss} zuerst die Kreistheilungsgleichungen behandelt\footnote{\gesperrt{Gauss}, Disquisitiones arithmeticae, Sectio VII.}.

Wir knüpfen endlich noch die für die Folge wichtige Bemerkung hier an, dass für einen Primzahlgrad die Begriffe der \gesperrt{Normalgleichung und der cyklischen Gleichung zu}sammenfallen. Denn jedenfalls ist eine cyklische Gleichung vom Primzahlgrad, da sie irreducibel ist, eine Normalgleichung. Und wenn umgekehrt der Grad $n$ einer Normalgleichung, der zugleich der Grad der Gruppe dieser Gleichung ist, eine Primzahl ist, und $\pi$ eine nicht identische Permutation dieser Gruppe, so ist der Grad von $\pi$, der ja ein Theiler von $n$ sein muss, gleich $n$, und die Gruppe der Gleichung ist $1, \pi, \pi^2 \ldots \pi^{n-1}$, also cyklisch.
